% Generated by scripts/build-paper.py — edit paper/src/survey.src.md
\section{Review method}\label{sec:review-method}

\subsection{Search}

The snapshot of 23 September 2026 ran 17 accepted arXiv API queries. They returned 225 hits and 197 unique records after de-duplication by unversioned arXiv ID (Appendix A). Queries combined the brand and interface vocabulary (\texttt{jev}, \texttt{typesafe}, “System One”, “typed decision”, “calibrated decisions”) with open-model names, survey terms and a date-bounded “decision model” query; arXiv \texttt{all:} searches metadata rather than full text, and that last query found a core study the brand-name queries missed \citep{arxiv2609_24574}. A same-day increment added 19 expansion queries; queries whose phrases the API parsed so loosely that they returned thousands of records were rejected by rule.

Later increments re-ran every accepted and expansion query after the daily arXiv announcement and screened only records not seen before. The increment of 24 September added four core studies. The increment of 6 October retrieved 710 hits (444 unique records), of which 112 had not been screened: 74 became core studies, six were added as background or related reviews, and 32 were excluded with a recorded reason. Each run also refreshed the metadata of every included record; seven included studies had posted revised versions, which were re-read and their records updated. The increment of 7 October, run once that day's listing had reached the API, retrieved 743 hits (458 unique records); ten of its 14 new records became core studies and four were excluded, and of three revised studies one had been substantially rewritten and was re-read. Every screening decision and reason is published with the data.

Repositories came from GitHub searches and from curated lists. For the lists we checked each entry against its own README or paper before inclusion: entries whose README did not mention Jev or a typed decision model were excluded, and general infrastructure or prior work listed for context was recorded as context only \citep{gh_omnijev_awesome_jev_gallery}.

\subsection{Eligibility}

A study is \textbf{core} if it evaluates TypeSafe Jev, a clearly identified Jev-like typed-decision implementation, or a system whose contribution materially depends on such decisions. \textbf{Peripheral} studies fall within reach of the scope but rest on contested, self-reported claims; they stay visible but are not pooled. \textbf{Background} references are adjacent methods, related reviews and earlier work the interface descends from, each with a stated role. We excluded name matches (Japanese encephalitis, JEPA, unrelated uses of “TypeSafe” or “System One”, other expansions of RLCD) and work without a typed-decision link.

\subsection{Reading depth and extraction}

We read the core and peripheral studies in full, checking key tables and limitation sections; background references were checked against metadata and abstracts. Code, weight and data links reported by the studies were checked to resolve, and releases that a paper announces but that could not be found are recorded as “claimed, not located”. Repository evidence comes from READMEs and metadata at pinned commits, with selected source files for a few. Nothing was executed and no experiment was re-run.

Each checkable statement became an evidence record. A record stores its subject, text, headline and source locator, and an evidence type: author-reported experiment, vendor documentation, vendor-reported result or community report. It also stores the metric, value and baseline; the sample size, model version and hardware, each with a reason when unknown; the \textbf{test level} (model test, system test, hybrid, simulation); and the \textbf{measurement scope} (single request, amortized per question, batch, end to end, simulation, author estimate, vendor claim). Each record links to the findings of §7 as \emph{supporting} or \emph{qualifying} evidence. Openness is recorded as six independent fields per study: code visible, weights available, data available, raw predictions available, recomputable, independently reproduced.

\subsection{Synthesis and its limits}

Tasks, metrics, binning and scopes are heterogeneous, so we pool nothing and rank nothing. Findings are stated at the strength their records allow, with qualifying records listed beside supporting ones. Studies that share authors are flagged as study families and read together: two edge and RAN studies by one group (which also wrote the network SoK), three studies by Rafe and Das, an agent architecture and an attack study by Wu and Lim, two studies by Yang and Zhao, a Turkish benchmark whose author also built a model it ranks, and a reproduction of Laya and a judge evaluation of CLM by one author \citep{arxiv2609_22753,arxiv2609_23136,arxiv2609_24052,arxiv2610_00213,arxiv2610_00346,arxiv2609_26532,arxiv2609_28613,arxiv2609_29283,arxiv2610_06354,arxiv2610_02293,arxiv2610_06744,arxiv2609_33843,arxiv2610_07177}. The review was carried out by one AI-assisted reviewer and is not a registered or PRISMA-compliant systematic review; §12 lists what this implies.
